\section{Path functionals on arbitrary subspaces}\label{branch:section}\label{sec:branches}

The exact pairs of Proposition~\ref{pair:exact} belong to block subspaces.
We now transfer their separation property to every infinite-dimensional
closed subspace of $E$.  The transfer uses a projection onto an initial
ordinal segment, followed, when possible, by a projection onto one color.
These projections preserve the path evaluations that we shall construct.

\begin{lemma}\label{branch:reduction}
Let $F$ be an infinite-dimensional closed subspace of $E$.  There are a
nonzero limit ordinal $\gamma\leq\omega_1$, an infinite-dimensional closed
subspace $H\subset P_{<\gamma}E$, a bounded linear map $U:H\to F$, and one
of the following two alternatives:
\begin{enumerate}
\item For some color $C_k$, we have $H\subset P_{C_k}E$ and
      $P_{C_k}P_{<\gamma}Uh=h$ for every $h\in H$.
\item We have $P_{<\gamma}Uh=h$ for every $h\in H$, and every restriction
      $P_{C_k}|_H$ is strictly singular.
\end{enumerate}
In both alternatives, $P_{<\beta}|_H$ is strictly singular for every
$\beta<\gamma$.  Moreover, if $P$ is an actual path functional supported
below $\gamma$, and supported on $C_k$ in the first alternative, then
\begin{equation}\label{branch:lift-identity}
                         P(Uh)=P(h)\qquad(h\in H).
\end{equation}
\end{lemma}

\begin{proof}
Choose $\gamma\leq\omega_1$ least such that $P_{<\gamma}|_F$ is not
strictly singular; it exists because $P_{<\omega_1}=I_E$.
It is nonzero.  It cannot be a successor $\beta+1$: on an
infinite-dimensional closed subspace where $P_{<\beta+1}$ is bounded
below, intersect with the kernel of $e_\beta^*$.  This intersection is
still infinite dimensional, and $P_{<\beta}=P_{<\beta+1}$ there,
contradicting minimality.

Choose an infinite-dimensional closed $F_0\subset F$ on which
$P_{<\gamma}$ is bounded below.  Its image $H_0=P_{<\gamma}F_0$ is
closed, and its inverse $U_0:H_0\to F_0$ is bounded.  Every earlier
projection is strictly singular on $H_0$.  Indeed, if $P_{<\beta}$ were
bounded below by $a>0$ on an infinite-dimensional closed $M\subset H_0$,
then
\[
 \|P_{<\beta}U_0h\|=\|P_{<\beta}h\|
       \geq a\|h\|\geq \frac{a}{\|U_0\|}\|U_0h\|
       \qquad(h\in M)
\]
would contradict the choice of $\gamma$.

If all the color projections are strictly singular on $H_0$, take
$H=H_0$ and $U=U_0$, giving the second alternative.  Otherwise choose
$k$ and an infinite-dimensional closed $H_1\subset H_0$ such that
$C=P_{C_k}$ is bounded below by $b>0$ on $H_1$.  Set
\[
 H=CH_1,\qquad U=U_0(C|_{H_1})^{-1}.
\]
The space $H$ is closed, and $CP_{<\gamma}U=I_H$.
To check earlier projections, suppose $P_{<\beta}$ is bounded below by
$a>0$ on an infinite-dimensional closed $M\subset H$.  On
$(C|_{H_1})^{-1}M$, commutation and contractivity of the projections give
\[
 \|P_{<\beta}z\|\geq\|CP_{<\beta}z\|
       =\|P_{<\beta}Cz\|\geq a\|Cz\|\geq ab\|z\|,
\]
contradicting strict singularity on $H_0$.

Finally, put $R=P_{<\gamma}$ in the second alternative and
$R=P_{C_k}P_{<\gamma}$ in the first.  Then $RU=I_H$.  Each finite
partial sum of a path with the stated support is fixed by composition
with $R$.  Passing to its pointwise limit on $E$ gives $P\circ R=P$,
and hence~\eqref{branch:lift-identity}.
\end{proof}

\begin{theorem}\label{branch:separation}
For every infinite-dimensional closed subspace $F\subset E$, there are
a single operation family, a number $\delta_F>0$, and actual path
functionals $P_b\in B_{E^*}$ of that family, indexed by
$b\in\{0,1\}^{\N}$, such that
\[
          \|P_b|_F-P_c|_F\|\geq\delta_F
          \qquad(b\ne c).
\]
\end{theorem}

\begin{proof}
Apply Lemma~\ref{branch:reduction}.  We use the pure $k$ family in its
first alternative and the mixed family in its second.  There are two
steps: approximate vectors in $H$ by successive finite blocks, and use
exact pairs in those blocks to label a binary tree.

\paragraph{Successive blocks and a small perturbation.}
Choose positive numbers $\varepsilon_i$ with
\[
           \sum_{i=1}^{\infty}\varepsilon_i\leq\frac1{32000}.
\]
We construct $h_i\in H$ and nonzero finite vectors $z_i\in E$ satisfying
\begin{equation}\label{branch:block-approximation}
 \|h_i\|=1,\qquad \|h_i-z_i\|<\varepsilon_i,
       \qquad \operatorname{supp}z_i\subset[0,\gamma),
\end{equation}
with the $z_i$ successive on $C_k$ in the pure case, and successive in
both ordinal and color ranges in the mixed case.

Here is the inductive step.  After $z_1,\ldots,z_{i-1}$ have been
chosen, let $\beta<\gamma$ lie strictly above their supports; this is
possible because $\gamma$ is a limit ordinal.  For $i=1$, take
$\beta=0$.  In the mixed case let $l$ be the largest color previously
used, with $l=-1$ at the first step.  Apply Lemma~\ref{ev:small} to
$P_{<\beta}|_H$, and in the mixed case also to the finitely many maps
$P_{C_j}|_H$, $0\leq j\leq l$.  We may therefore choose a unit vector
$h_i\in H$ for which
\[
 \|P_{<\beta}h_i\|<\varepsilon_i/3,
 \qquad
 \sum_{j=0}^{l}\|P_{C_j}h_i\|<\varepsilon_i/3
       \quad\hbox{in the mixed case}.
\]
The second sum is zero when $l=-1$.
In the pure case remove the initial ordinal segment from $h_i$; in the
mixed case remove also all colors at most $l$.  The resulting vector
differs from $h_i$ by less than $2\varepsilon_i/3$, because the
remaining ordinal projection is contractive.  It is supported below
$\gamma$ and in the required ordinal and color tails.  A finite-support
approximation, cropped to those same sets, gives $z_i$ satisfying
\eqref{branch:block-approximation}.  In the pure case the approximation
is additionally cropped to $C_k$.

The ordinal interval projections show that $(z_i)$ is a basic sequence
with prefix projections of norm at most one.  Since $\|z_i\|\geq1/2$,
its coefficients satisfy
\[
      |a_i|\leq4\Big\|\sum_j a_jz_j\Big\|
\]
on finite sums.  Consequently the map defined there by $Vz_i=h_i$
extends to a bounded map $V:[z_i]\to H$ with
\begin{equation}\label{branch:perturbation}
 \|V-I\|\leq4\sum_i\varepsilon_i\leq\eta,
          \qquad \eta:=\frac1{8000},
          \qquad \|V\|\leq1+\eta.
\end{equation}
Here $I:[z_i]\to E$ is the inclusion, so the first norm is taken for
maps into $E$; the image of $V$ lies in $H$ because $H$ is closed.

\paragraph{A binary tree of exact pairs.}
As in the James tree construction~\cite{James1974}, branches will
separate dual functionals. Here the nodes carry the prescribed-weight
pairs constructed in Proposition~\ref{pair:exact}.
List the finite binary sequences in breadth-first order, ordering each
finite level lexicographically.  For each node $s$ we shall choose a
finite block vector $y_s\in\operatorname{span}\{z_i:i\in\N\}$, a
Type~I functional $\phi_s$, its declared weight $J_s$, and its integer
metadata $p_s$.  Apply Proposition~\ref{pair:exact} successively in this
order, always beyond every block used at earlier nodes.  It gives
\begin{equation}\label{branch:node-pair}
             \|y_s\|\leq500,
             \qquad \phi_s(y_s)\geq\frac12.
\end{equation}
The support of $\phi_s$ lies in the ordinal block range used for $y_s$,
and in its color range as well in the mixed case.  Thus the supporting
ranges for different nodes are disjoint and successive in the chosen
enumeration.  In particular,
\begin{equation}\label{branch:cross-evaluations}
                       \phi_s(y_t)=0\qquad(s\ne t).
\end{equation}

We now specify the weights and metadata, so that every branch is an
actual path.  Give the root any positive weight.  Once $\phi_s$ is
chosen, choose $p_s$ greater than the metadata at its predecessor, and
at least $J_s$ and the maximum of $\rho$ on the finite union of supports
along the branch from the root to $s$.  At the root the predecessor
condition is omitted.  Form the prescribed order-position code of this
whole prefix, and let both children of $s$ have weight equal to its
$\sigma$-value.  Their weight is therefore already known when they are
processed.  Proposition~\ref{pair:exact} allows any target weight in any tail,
so every step is possible.  The siblings have the same weight; this
causes no conflict, since they never belong to the same path.  The
disjoint-active-weight requirement for a Type~II combination places no
additional restriction on the individual paths constructed here.

For $b\in\{0,1\}^{\N}$, let
\[
                  P_b=\sum_{n=0}^{\infty}\phi_{b|n},
\]
where $b|0$ is the root.  Succession and the chosen codes make this an
actual path in the selected family.  Its finite partial sums have norm
at most one, so it defines $P_b\in B_{E^*}$.  Every piece is supported
below $\gamma$, and on $C_k$ in the pure case.  Thus
\eqref{branch:lift-identity} applies to every $P_b$.

If $b\ne c$, choose a node $s$ on $b$ beyond their first divergence.
Equations~\eqref{branch:node-pair} and
\eqref{branch:cross-evaluations} give
\[
                       P_b(y_s)\geq\frac12,
                       \qquad P_c(y_s)=0.
\]
Test their restrictions to $F$ at $UVy_s$.  By the lift identity and
\eqref{branch:perturbation},
\begin{align*}
 (P_b-P_c)(UVy_s)
   &=(P_b-P_c)(Vy_s)\\
   &\geq\frac12-2\eta\|y_s\|
    \geq\frac38,
 \qquad
 \|UVy_s\|\leq500\|U\|(1+\eta).
\end{align*}
The conclusion follows with
\[
                \delta_F=\frac{3}{8\cdot500\|U\|(1+\eta)}>0.
\]
\end{proof}
